\chapter{Trabalhos Relacionados}
\label{cap:trabalhos-relacionados}
O desenvolvimento de Sistemas de Aquisição de Dados (SAD) e telemetria embarcada para Veículos Aéreos Não Tripulados (VANTs) de pequeno porte, especialmente no cenário de validação de projetos aeroespaciais (como na competição SAE Brasil AeroDesign), é um tema recorrente na literatura acadêmica. A necessidade empírica de validar modelos teóricos de aerodinâmica, desempenho e cargas estruturais impulsiona a criação de soluções customizadas que atendam a rigorosas restrições de SWaP (Size, Weight, and Power).

Trabalhos como o de Barbosa \cite{barbosa2022} e Bueno \cite{bueno2015} ilustram a abordagem clássica de engenharias universitárias para este problema. Barbosa \cite{barbosa2022} desenvolveu um sistema de telemetria focado no monitoramento em tempo real de variáveis de voo para projetos de aeromodelismo, com foco em eletrônica capaz de prover suporte computacional para testes. Semelhantemente, Bueno \cite{bueno2015} investigou a arquitetura de software embarcado de tempo real para instrumentação aeroelástica, frequentemente empregando módulos de comunicação da família XBee/ZigBee (2.4 GHz). Embora eficazes para coleta primária, essas soluções comumente adotam topologias fechadas (ponto-a-ponto monolítico) e protocolos de empacotamento em texto plano ou matrizes de bytes proprietárias, limitando severamente a escalabilidade caso o projeto demande integração de estações meteorológicas externas ou múltiplos VANTs simultâneos.

No que tange à camada física e de enlace, a adoção da modulação LoRa (Long Range) tem ganhado destaque acadêmico na instrumentação de VANTs por sua resiliência e alcance superior operando em frequências sub-GHz (915 MHz). Contudo, a literatura aponta o limitante crítico do LoRa: a restrita capacidade de canal, descrita pelo Teorema de Shannon-Hartley. Projetos de baixo custo geralmente contornam essa limitação enviando strings genéricas. Já os frameworks padronizados na indústria, como o protocolo MAVLink (adotado pelo ArduPilot e PX4), são comumente aplicados sobre links de telemetria dedicados de maior banda ou delegados a pesadas controladoras de voo prontas (Pixhawk), conforme analisado em estudos recentes de tráfego de rede para drones \cite{belwafi2022}. A geração de mensagens MAVLink diretamente em microcontroladores de borda (edge computing) em redes distribuídas de baixa banda continua sendo um desafio tecnológico no meio acadêmico de baixo orçamento.

Neste contexto, o presente trabalho contextualiza-se na mesma demanda de instrumentação aeronáutica acessível, porém diverge e inova substancialmente frente às implementações acadêmicas vigentes em três pilares metodológicos:

- Arquitetura Distribuída e Middleware: Ao contrário de módulos isolados que conectam a aeronave de forma unilateral ao visualizador no solo, o presente projeto implementa um modelo distribuído através do Shoulder (middleware). Essa camada centraliza, valida e formata os pacotes provenientes de nós completamente independentes (sistema da aeronave, obstáculos e clima terrestre), atuando como um barramento unificador de eventos, distanciando-se das arquiteturas monolíticas triviais.

- Empacotamento MAVLink com Otimização por Quantização: Diferentemente das soluções universitárias que evitam headers padronizados para economizar processamento, este trabalho aplica o dialeto MAVLink (v1/v2) como contrato de rede base, preservando a interoperabilidade. A limitação física do LoRa é solucionada de forma analítica: o domínio executa a quantização matemática em ponta, convertendo variáveis nativas em ponto flutuante ($float32$) para formatos inteiros escalonados ($int16\_t$ e $int32\_t$). Esta mitigação drástica de payload raramente é abordada em soluções equivalentes baseadas na abstração excessiva do ecossistema Arduino.

- Rigor de Engenharia de Software Orientado a Domínio: Na literatura de instrumentação maker, os firmwares são majoritariamente codificados em formato bare-metal procedural (Super Loops). Este trabalho opta pela implementação sob o FreeRTOS através de Arquitetura Hexagonal (Ports and Adapters). Essa separação técnica desvincula a regra de negócios aeroespacial dos periféricos elétricos, viabilizando ensaios de injeção de erros (Hardware-in-the-loop - HIL) antes da manufatura final da PCB, elevando o grau de determinismo e modularidade do sistema.

Portanto, o sistema abordado supera o escopo tradicional de data loggers isolados, assumindo a estrutura de um ecossistema de IoT aeronáutico altamente determinístico e focado no balanço ótimo entre restrição física, conformidade de rede aeroespacial e escalabilidade de engenharia de ensaios em voo.
